\subsection{The existence and uniqueness theorem}

DEFINITION 2. --- Let G be a locally compact group. A nonzero positive measure on G that is left (resp. right) invariant is called a left (resp. right) Haar measure on G.

THEOREM 1. --- On every locally compact group, there exists a left (resp. right) Haar measure, and, up to a constant factor, there exists only one.

\begin{enumerate}
\def\labelenumi{\Alph{enumi})}
\tightlist
\item
  Existence. --- Set \(\mathcal{K}(\mathrm{G}) = \mathcal{K}\) , \(\mathcal{K}_{+}(\mathrm{G}) = \mathcal{K}_{+}\) , \(\mathcal{K}_{+}^{*} = \mathcal{K}_{+} - \{0\}\) . If C is a compact subset of G, we denote by \(\mathcal{K}_{+}^{*}(\mathrm{C})\) the set of \(f \in \mathcal{K}_{+}^{*}\) with support in C. For \(f \in \mathcal{K}\) and \(g \in \mathcal{K}_{+}^{*}\) , there exist numbers \(c_{1}, \ldots, c_{n} \geqslant 0\) and elements \(s_{1}, \ldots, s_{n}\) of G such that \(f \leqslant \sum_{i=1}^{n} c_{i} \gamma(s_{i}) g\) : for, there exists a nonempty open set U in G such that \(\inf_{s \in \mathrm{U}} g(s) > 0\) , and the support of \(f\) can be covered by a finite number of left-translates of U. Let \((f : g)\) be the infimum of the numbers \(\sum_{i=1}^{n} c_{i}\) for all systems \((c_{1}, \ldots, c_{n}, s_{1}, \ldots, s_{n})\) of numbers \(\geqslant 0\) and elements of G such that \(f \leqslant \sum_{i=1}^{n} c_{i} \gamma(s_{i}) g\) . Then:
\end{enumerate}

\begin{enumerate}
\def\labelenumi{(\roman{enumi})}
\item
  \(\left(\pmb{\gamma}(s)f:g\right) = (f:g)\) for \(f\in \mathcal{K}\) , \(g\in \mathcal{K}_{+}^{*}\) , \(s\in G\) ;
\item
  \((\lambda f:g) = \lambda (f:g)\) for \(f\in \mathcal{K}\) \(,g\in \mathcal{K}_{+}^{*},\lambda \geqslant 0\)
\item
  \(\left((f + f^{\prime}):g\right)\leqslant (f:g) + (f^{\prime}:g)\) for \(f\in \mathcal{K},f^{\prime}\in \dot{\mathcal{K}},g\in \mathcal{K}_{+}^{*};\)
\item
  \((f:g)\geqslant (\sup f) / (\sup g)\)
\end{enumerate}

\[
f \in \mathscr {K}, g \in \mathscr {K} _ {+} ^ {*}
\]

\begin{enumerate}
\def\labelenumi{(\alph{enumi})}
\setcounter{enumi}{21}
\tightlist
\item
  \((f:h)\leqslant (f:g)(g:h)\) for \(f\in \mathcal{K}\) \(,g\in \mathcal{K}_{+}^{*}\) \(h\in \mathcal{K}_{+}^{*};\)
\end{enumerate}

\begin{enumerate}
\def\labelenumi{(\roman{enumi})}
\setcounter{enumi}{5}
\item
  \(0 < \frac{1}{(f_0:f)}\leqslant \frac{(f:g)}{(f_0:g)}\leqslant (f:f_0)\) for \(f,f_0,g\) in \(\mathcal{K}_+^*\) ;
\item
  let \(f, f', h\) be in \(\mathcal{K}_+\) with \(h(s) \geqslant 1\) on the support of \(f + f'\) , and let \(\varepsilon > 0\) ; there exists a compact neighborhood \(V\) of \(e\) such that, for every \(g \in \mathcal{K}_+^*(V)\) ,
\end{enumerate}

\[
(f: g) + (f ^ {\prime}: g) \leqslant ((f + f ^ {\prime}): g) + \varepsilon (h: g).
\]

The properties (i), (ii), (iii) are obvious. Let \(f \in \mathcal{K}\) , \(g \in \mathcal{K}_+^*\) ; if \(f \leqslant \sum_{i=1}^{n} c_i \gamma(s_i)g\) with the \(c_i \geqslant 0\) , then \(\sup f \leqslant \sum_{i=1}^{n} c_i g(s_i^{-1}s)\) for some \(s \in \mathbf{G}\) , therefore \(\sup f \leqslant \left( \sum_{i=1}^{n} c_i \right) \sup g\) , whence (iv). Let us now prove (v); let \(f \in \mathcal{K}\) , \(g, h\) in \(\mathcal{K}_+^*\) ; if \(f \leqslant \sum_{i=1}^{n} c_i \boldsymbol{\gamma}(s_i) g\) and \(g \leqslant \sum_{j=1}^{p} d_j \boldsymbol{\gamma}(t_j) h\) ( \(c_i \geqslant 0\) , \(d_j \geqslant 0\) , \(s_i, t_j\) in G), then \(f \leqslant \sum_{i,j} c_i d_j \boldsymbol{\gamma}(s_i t_j) h\) , therefore \((f: h) \leqslant \sum_{i,j} c_i d_j = \left( \sum_{i} c_i \right) \left( \sum_{j} d_j \right)\) ; thus \((f: h) \leqslant (f: g)(g: h)\) . Applying (v) to \(f_0, f, g\) on the one hand and to \(f, f_0, g\) on the other, one obtains (vi). Finally, let \(f, f', h\) be in \(\mathcal{K}_+\) with \(h(s) \geqslant 1\) on the support of \(f + f'\) , and let \(\varepsilon > 0\) . Set \(\mathbf{F} = f + f' + \frac{1}{2} \varepsilon h\) ; the functions \(\varphi, \varphi'\) , that coincide respectively with \(f / F\) and \(f'/F\) on the support of \(f + f'\) and are zero outside it, belong to \(\mathcal{K}_+\) ; for every \(\eta > 0\) , there exists a compact neighborhood V of \(e\) such that \(|\varphi(s) - \varphi(t)| \leqslant \eta\) and \(|\varphi'(s) - \varphi'(t)| \leqslant \eta\) for \(s^{-1} t \in V\) . Then let \(g \in \mathcal{K}_+^*(V)\) ; for every \(s \in G\) ,

\[
\varphi \cdot \boldsymbol {\gamma} (s) g \leqslant (\varphi (s) + \eta) \cdot \boldsymbol {\gamma} (s) g;
\]

for, this is obvious at the points where \(\pmb{\gamma}(s)g\) is zero, therefore outside of \(s\mathbf{V}\) ; and in \(s\mathbf{V}\) , \(\varphi \leqslant \varphi(s) + \eta\) ; similarly, \(\varphi' \cdot \pmb{\gamma}(s)g \leqslant (\varphi'(s) + \eta) \cdot \pmb{\gamma}(s)g\) . This stated, let \(c_1, \ldots, c_n\) be numbers \(\geqslant 0\) and \(s_1, \ldots, s_n\) elements of \(\mathbf{G}\) such that \(\mathbf{F} \leqslant \sum_{i=1}^{n} c_i \pmb{\gamma}(s_i)g\) ; then

\[
f = \varphi \mathrm{F} \leqslant \sum_ {i = 1} ^ {n} c _ {i} \varphi \cdot \boldsymbol {\gamma} (s _ {i}) g \leqslant \sum_ {i = 1} ^ {n} c _ {i} (\varphi (s _ {i}) + \eta) \cdot \boldsymbol {\gamma} (s _ {i}) g
\]

and similarly for \(f'\) ; consequently

\[
(f: g) + (f ^ {\prime}: g) \leqslant \sum_ {i = 1} ^ {n} c _ {i} (\varphi (s _ {i}) + \varphi^ {\prime} (s _ {i}) + 2 \eta) \leqslant (1 + 2 \eta) \sum_ {i = 1} ^ {n} c _ {i}
\]

since \(\varphi +\varphi^{\prime}\leqslant 1\) . Applying the definition of F, then (ii), (iii) and (v), one infers that

\[
\begin{array}{r l} & {(f: g) + (f ^ {\prime}: g) \leqslant (1 + 2 \eta) (\mathrm{F}: g) \leqslant} \\ & {\quad (1 + 2 \eta) \big [ \big ((f + f ^ {\prime}): g \big) + \frac {1}{2} \varepsilon (h: g) \big ] \leqslant} \\ & {\quad \big ((f + f ^ {\prime}): g \big) + \frac {1}{2} \varepsilon (h: g) + 2 \eta \big ((f + f ^ {\prime}): h \big) (h: g) + \varepsilon \eta (h: g)} \end{array}
\]

and, if \(\eta\) has been chosen so that \(\eta\big[2\big((f + f'):h\big) + \varepsilon\big] \leqslant \frac{1}{2}\varepsilon\) , one obtains (vii).

As V runs over the set of compact neighborhoods of \(e\) , the \(\mathcal{K}_{+}^{*}(V)\) form a base of a filter \(\mathfrak{B}\) on \(\mathcal{K}_{+}^{*}\) . Let \(\mathfrak{F}\) be an ultrafilter on \(\mathcal{K}_{+}^{*}\) finer than \(\mathfrak{B}\) . On the other hand, let us fix \(f_{0} \in \mathcal{K}_{+}^{*}\) and let us set, for \(f \in \mathcal{K}_{+}^{*}\) and \(g \in \mathcal{K}_{+}^{*}\) ,

\[
\mathrm{I} _ {g} (f) = \frac {(f : g)}{(f _ {0} : g)}.
\]

By (vi), \(\lim_{g,\mathfrak{F}}\mathrm{I}_g(f) = \mathrm{I}(f)\) exists in the compact space \([1 / (f_0:f),(f:f_0)]\) . By (iii), \(\mathrm{I}(f + f')\leqslant \mathrm{I}(f) + \mathrm{I}(f')\) . By (vii), \(\mathrm{I}(f) + \mathrm{I}(f')\leqslant \mathrm{I}(f + f') + \varepsilon \mathrm{I}(h)\) for all \(\varepsilon >0\) provided \(h\) is \(\geqslant 1\) on the support of \(f + f'\) ; it follows that \(\mathrm{I}(f + f') = \mathrm{I}(f) + \mathrm{I}(f')\) . By Ch. II, §2, No.~1, Prop. 2, I is extendible to a linear form on \(\mathcal{K}\) ; this linear form is a nonzero positive measure on G, left-invariant by (i); this is the sought-for left Haar measure. Passing to the opposite group, one deduces from this the existence of a right Haar measure.

\begin{enumerate}
\def\labelenumi{\Alph{enumi})}
\setcounter{enumi}{1}
\tightlist
\item
  Uniqueness. --- Let \(\mu\) be a left Haar measure, \(\nu\) a right Haar measure. Then \(\check{\nu}\) is a left Haar measure. We are going to show that \(\mu\) and \(\check{\nu}\) are proportional. This will prove that any two left Haar measures are indeed proportional.
\end{enumerate}

Let \(f \in \mathcal{K}\) be such that \(\mu(f) \neq 0\) . By Lemma 1, the function \(D_f\) defined on \(G\) by the formula

\[
\mathrm{D} _ {f} (s) = \mu (f) ^ {- 1} \int f (t ^ {- 1} s) d \nu (t)\tag{16}
\]

is continuous on G. Let \(g \in \mathcal{K}\) . The function \((s, t) \mapsto f(s)g(ts)\) is continuous with compact support in G × G. By Ch. III, §4, No.~1, Th. 2,

\[
\begin{array}{r l} & {\mu (f) \nu (g) = \Big (\int f (s) d \mu (s) \Big) \Big (\int g (t) d \nu (t) \Big)} \\ & {\qquad = \int d \mu (s) \int f (s) g (t s) d \nu (t) = \int d \nu (t) \int f (s) g (t s) d \mu (s)} \\ & {\qquad = \int d \nu (t) \int f (t ^ {- 1} s) g (s) d \mu (s)} \\ & {\qquad = \int g (s) \Big [ \int f (t ^ {- 1} s) d \nu (t) \Big ] d \mu (s) = \mu \big (g \cdot \mu (f) D _ {f} \big),} \end{array}\tag{17}
\]

whence

\[
\nu (g) = \mu (\mathrm{D} _ {f} \cdot g).\tag{18}
\]

This proves, first, that \(D_{f}\) does not depend on f.~For, if \(f' \in H\) is such that \(\mu(f') \neq 0\) , then \(D_{f} \cdot \mu = D_{f'} \cdot \mu\) , therefore \(D_{f} = D_{f'}\) locally almost everywhere for \(\mu\) , hence everywhere, since \(\mathbf{D}_f\) and \(\mathbf{D}_{f'}\) are continuous and the support of \(\mu\) is \(\mathbf{G}\) . We may therefore write \(\mathrm{D}_f = \mathrm{D}\) . The formula (16) gives

\[
\mu (f) \mathrm{D} (e) = \stackrel {\vee} {\nu} (f).\tag{19}
\]

The formula (19) may be extended by linearity to the functions \(f \in \mathcal{K}\) such that \(\mu(f) = 0\) . We have \(D(e) \neq 0\) since \(\check{\nu} \neq 0\) . This indeed establishes the proportionality of \(\mu\) and \(\check{\nu}\) .

COROLLARY. --- Every left-invariant (resp. right-invariant) measure on G is proportional to a left (resp. right) Haar measure.

Examples. --- 1) On the additive group R, the Lebesgue measure dx is a Haar measure (Ch. III, §1, No.~3, Example).

\begin{enumerate}
\def\labelenumi{\arabic{enumi})}
\setcounter{enumi}{1}
\tightlist
\item
  For every function \(f \in \mathcal{K}(\mathbf{R}_+^*)\) , we have (FRV, II, §1, formula (12))
\end{enumerate}

\[
\int_ {0} ^ {+ \infty} \frac {f (x)}{x} d x = \int_ {0} ^ {+ \infty} \frac {f (t x)}{t x} t d x = \int_ {0} ^ {+ \infty} \frac {f (t x)}{x} d x
\]

for all t \textgreater{} 0; the measure \(x^{-1} dx\) is thus a Haar measure on the multiplicative group \(R_{+}^{*}\) .

\begin{enumerate}
\def\labelenumi{\arabic{enumi})}
\setcounter{enumi}{2}
\tightlist
\item
  Let us take for G the torus \(\mathbf{T} = \mathbf{R} / \mathbf{Z}\) . Let \(\varphi\) be the canonical mapping of \(\mathbf{R}\) onto \(\mathbf{T}\) . For \(f \in \mathcal{K}(\mathbf{T})\) , the function \(f \circ \varphi\) is continuous and periodic with period 1 on \(\mathbf{R}\) , and the integral
\end{enumerate}

\[
\mathrm{I} (f) = \int_ {a} ^ {a + 1} f (\varphi (x)) d x
\]

is independent of the choice of \(a \in \mathbf{R}\) ; it is immediate that it is invariant under translation; it therefore defines a Haar measure on \(\mathbf{T}\) . By transport of structure, one deduces from this that \(I(f) = \int_{a}^{a + 1} f(e^{2\pi it}) dt\) is a Haar measure on the multiplicative group \(\mathbf{U}\) of complex numbers of absolute value 1 (GT, VIII, §2, No.~1).

PROPOSITION 2. --- Let G be a locally compact group, \(\mu\) a left or right Haar measure on G. For G to be discrete, it is necessary and sufficient that \(\mu(\{e\}) > 0\) . For G to be compact, it is necessary and sufficient that \(\mu^{*}(G) < +\infty\) .

The conditions are obviously necessary. Let us show their sufficiency. Let V be a compact neighborhood of e. If \(\mu(\{e\}) > 0\) , then V is a finite set since \(\mu(V) < +\infty\) ; since G is Hausdorff, it is therefore discrete. Suppose \(\mu^{*}(G)<+\infty\) and \(\mu\) is, for example, left-invariant. Consider the set \(\mathscr{E}\) of finite subsets \(\{s_{1},\ldots,s_{n}\}\) of G such that \(s_{i}V\cap s_{j}V=\varnothing\) for \(i\neq j\) ; then

\[
n \mu (\mathrm{V}) = \mu (s _ {1} \mathrm{V} \cup \dots \cup s _ {n} \mathrm{V}) \leqslant \mu^ {*} (\mathrm{G}),
\]

therefore \(n \leqslant \mu^{*}(\mathrm{G})/\mu(\mathrm{V})\) . We may therefore choose in E a maximal element \(\{s_{1},\ldots,s_{n}\}\) . Then, for every \(s \in G\) , there exists an i such that \(sV \cap s_{i}V \neq \varnothing\) , hence such that \(s \in s_{i}VV^{-1}\) . Thus G is the union of the compact sets \(s_{i}VV^{-1}\) , hence is compact.

